%%%%%%%%%%%%%%%%%%%%%%%%%%%%%%%%%%%%%%%%%%%%%%%%%%%%%%%%%%%%%%%%%%%%%%%%%%%%%%%%
%% RGAIG+ STANDARDS, DEVICE BENCHMARKS, AND CLINICAL VALIDATION TABLES
%% Covers: Regulatory standards, device benchmarks, safety, clinical validation,
%% consumer vs clinical EEG comparison, human factors, standards-layer mapping
%%%%%%%%%%%%%%%%%%%%%%%%%%%%%%%%%%%%%%%%%%%%%%%%%%%%%%%%%%%%%%%%%%%%%%%%%%%%%%%%
\normalsize\normalfont\color{black}

%% ── Table 1: Regulatory Standards Alignment Matrix ──
Table~\ref{tab:rgaig_standards_alignment} presents the fifteen regulatory standards integrated within the RGAIG+ framework, organised by governance domain.

\needspace{4\baselineskip}
\begingroup\singlespacing\tablefontsize\tablestripes
\begin{xltabular}{\textwidth}{@{} p{2.2cm} p{2.2cm} X p{1.8cm} @{}}
\caption{RGAIG+ Regulatory Standards Alignment Matrix}
\label{tab:rgaig_standards_alignment}\\
\toprule
\thead{Standard} & \thead{Domain} & \thead{Purpose and Scope} & \thead{Region} \\
\midrule
\endfirsthead
\endhead
\bottomrule
\endlastfoot
\multicolumn{4}{@{}p{\dimexpr\textwidth-2\tabcolsep}@{}}{\textbf{\color{ggunavy}Medical Device Safety}} \\
\midrule
IEC 60601                & Device safety    & Electrical safety and essential performance for medical electrical equipment; applies to consumer EEG wearables providing diagnostic insights & Global \\
\addlinespace
ISO 14971                & Risk management  & Application of risk management to medical devices; hazard identification, risk estimation, risk evaluation, and residual risk acceptability & Global \\
\addlinespace
ISO 13485                & Quality system   & Quality management system requirements for medical device design, development, production, and post-market surveillance & Global \\
\addlinespace
\midrule
\multicolumn{4}{@{}p{\dimexpr\textwidth-2\tabcolsep}@{}}{\textbf{\color{ggunavy}AI Governance}} \\
\midrule
ISO/IEC 42001            & AI governance    & Management system standard for organisations developing, providing, or using AI systems; risk-based approach to AI lifecycle governance & Global \\
\addlinespace
NIST AI RMF              & AI risk          & Voluntary framework for managing AI risks: Map, Measure, Manage, Govern functions; applicable across AI lifecycle & US \\
\addlinespace
OECD AI Principles       & Responsible AI   & Five principles for trustworthy AI: inclusive growth, human-centred values, transparency, robustness, accountability & OECD (38 countries) \\
\addlinespace
EU AI Act                & AI regulation    & Risk-based classification of AI systems (unacceptable, high, limited, minimal risk); mandatory requirements for high-risk AI in healthcare & EU \\
\addlinespace
\midrule
\multicolumn{4}{@{}p{\dimexpr\textwidth-2\tabcolsep}@{}}{\textbf{\color{ggunavy}Software as a Medical Device (SaMD)}} \\
\midrule
IMDRF SaMD Guidance      & Software device  & International framework classifying software intended for medical purposes; risk categorisation and clinical evaluation requirements & Global \\
\addlinespace
FDA AI/ML SaMD           & AI medical SW    & Regulatory guidance for AI/ML-based Software as a Medical Device; predetermined change control plan for adaptive algorithms & US \\
\addlinespace
\midrule
\multicolumn{4}{@{}p{\dimexpr\textwidth-2\tabcolsep}@{}}{\textbf{\color{ggunavy}Data Privacy and Security}} \\
\midrule
HIPAA                    & Health data      & Privacy and security rules for protected health information (PHI); applicable to EEG neurological data in US healthcare contexts & US \\
\addlinespace
GDPR                     & Personal data    & Data protection regulation; lawful basis, data minimisation, right to explanation, data portability; applies to EU subjects' EEG data & EU \\
\addlinespace
ISO/IEC 27001            & Cybersecurity    & Information security management system; risk-based controls for confidentiality, integrity, and availability of health data & Global \\
\addlinespace
\midrule
\multicolumn{4}{@{}p{\dimexpr\textwidth-2\tabcolsep}@{}}{\textbf{\color{ggunavy}Human Factors and Usability}} \\
\midrule
IEC 62366-1              & Usability        & Application of usability engineering to medical devices; use-related risk analysis, formative and summative evaluation & Global \\
\addlinespace
ISO 9241-210             & User experience  & Human-centred design for interactive systems; iterative design process with user involvement and evaluation & Global \\
\addlinespace
\midrule
\multicolumn{4}{@{}p{\dimexpr\textwidth-2\tabcolsep}@{}}{\textbf{\color{ggunavy}AI Safety and Monitoring}} \\
\midrule
WHO AI for Health        & Health AI        & Guidance on ethics and governance of AI for health; six guiding principles including transparency, accountability, and equity & Global \\
\addlinespace
IEEE 7000 Series         & Ethical AI       & Standards for ethical considerations in system design; addresses well-being, transparency, and accountability in AI systems & Global \\
\addlinespace
\midrule
\multicolumn{4}{@{}p{\dimexpr\textwidth-2\tabcolsep}@{}}{\textit{RGAIG+ integrates these 15 standards across its 9-layer architecture. Standards mapping to specific framework layers is provided in Table~\ref{tab:rgaig_standards_layer_map}. Compliance evidence tracked through the 525-module quality taxonomy (Section~3.8) and governance audit trail.}} \\
\end{xltabular}
\endgroup

\vspace{1.5em}

%% ── Table 2: Standards-to-Framework Layer Mapping ──
Table~\ref{tab:rgaig_standards_layer_map} maps each RGAIG+ framework layer to its primary regulatory standard and governance control.

\needspace{4\baselineskip}
\begingroup\singlespacing\tablefontsize\tablestripes
\begin{xltabular}{\textwidth}{@{} p{2.8cm} p{3cm} X @{}}
\caption{RGAIG+ Standards-to-Framework Layer Mapping}
\label{tab:rgaig_standards_layer_map}\\
\toprule
\thead{Framework Layer} & \thead{Primary Standard(s)} & \thead{Governance Control} \\
\midrule
\endfirsthead
\endhead
\bottomrule
\endlastfoot
1. Device Layer              & IEC 60601, ISO 13485          & Electrical safety certification; quality management for wearable EEG hardware; electrode biocompatibility \\
\addlinespace
2. Signal Processing         & ISO 14971                     & Risk management for signal artefact removal; error propagation analysis; signal quality thresholds (SNR $\geq$ 10~dB) \\
\addlinespace
3. Diagnostic AI             & ISO/IEC 42001, FDA SaMD       & AI model governance; predetermined change control plan; performance monitoring against clinical benchmarks \\
\addlinespace
4. Generative AI             & EU AI Act, OECD Principles    & High-risk AI classification; transparency requirements; hallucination control; explanation faithfulness \\
\addlinespace
5. Agentic Orchestration     & NIST AI RMF                   & Risk mapping for autonomous pipeline decisions; human-in-the-loop escalation; orchestration audit logging \\
\addlinespace
6. Consumer Interface        & IEC 62366-1, ISO 9241-210     & Usability engineering; human-centred design; use-related risk analysis; accessibility compliance \\
\addlinespace
7. Clinical Oversight        & IMDRF SaMD, WHO AI Health     & Clinical validation requirements; clinician-in-the-loop for high-risk outputs; informed consent mechanisms \\
\addlinespace
8. Monitoring                & ISO/IEC 27001, IEEE 7000      & Information security for monitoring data; ethical considerations in continuous surveillance; drift detection protocols \\
\addlinespace
9. Governance                & HIPAA, GDPR, ISO/IEC 42001    & Data privacy compliance; cross-jurisdictional governance; audit trail maintenance; regulatory reporting \\
\addlinespace
\midrule
\multicolumn{3}{@{}p{\dimexpr\textwidth-2\tabcolsep}@{}}{\textit{Each layer implements controls from its primary standard(s). Cross-cutting standards (GDPR, ISO 27001) apply across all layers. Compliance verified through 525-module quality taxonomy pass rates.}} \\
\end{xltabular}
\endgroup

\vspace{1.5em}

%% ── Table 3: Wearable Device Benchmark Framework ──
Table~\ref{tab:rgaig_device_benchmark} reports the eight-dimension benchmark framework for evaluating consumer wearable EEG devices.

\needspace{4\baselineskip}
\begingroup\singlespacing\tablefontsize\tablestripes
\begin{xltabular}{\textwidth}{@{} p{2.2cm} p{2.2cm} X p{2cm} @{}}
\caption{RGAIG+ Wearable EEG Device Benchmark Framework}
\label{tab:rgaig_device_benchmark}\\
\toprule
\thead{Dimension} & \thead{Benchmark Focus} & \thead{Description} & \thead{Target Metric} \\
\midrule
\endfirsthead
\endhead
\bottomrule
\endlastfoot
Signal quality       & EEG accuracy           & Signal-to-noise ratio of acquired EEG; correlation with clinical-grade systems; frequency band preservation (delta through gamma) & SNR $\geq$ 10~dB \\
\addlinespace
Device reliability   & Hardware stability      & Continuous uptime without signal dropout; electrode contact consistency over multi-hour recording sessions & Uptime $\geq$ 95\% \\
\addlinespace
Usability            & Ease of use            & Setup time from unboxing to first recording; self-application success rate; System Usability Scale score & SUS $\geq$ 70 \\
\addlinespace
Comfort              & Wearable ergonomics    & User comfort during extended wear ($>$ 2 hours); pressure distribution; weight; thermal comfort & Comfort $\geq$ 4/5 \\
\addlinespace
Battery performance  & Operational time       & Continuous recording duration per charge; standby time; charge-to-full duration & $\geq$ 6 hours \\
\addlinespace
Connectivity         & Data transmission      & Bluetooth/Wi-Fi latency from device to processing system; data loss rate during transmission; reconnection time & Latency $<$ 50~ms \\
\addlinespace
Calibration          & Signal consistency     & Electrode impedance stability over time; calibration drift between sessions; inter-session reproducibility & Drift $<$ 5\% \\
\addlinespace
Safety               & Electrical safety      & Current leakage limits per IEC 60601; electrode biocompatibility; skin irritation risk assessment & IEC 60601 compliant \\
\addlinespace
\midrule
\multicolumn{4}{@{}p{\dimexpr\textwidth-2\tabcolsep}@{}}{\textit{Benchmarks derived from IEC 60601, IEC 62366-1, and published consumer EEG validation studies. Target metrics represent minimum thresholds for consumer deployment readiness.}} \\
\end{xltabular}
\endgroup

\vspace{1.5em}

%% ── Table 4: Consumer vs Clinical EEG System Comparison ──
Table~\ref{tab:rgaig_consumer_vs_clinical} compares consumer wearable EEG and clinical EEG systems across eight performance and usability metrics.

\needspace{4\baselineskip}
\begingroup\singlespacing\tablefontsize\tablestripes
\begin{xltabular}{\textwidth}{@{} p{2cm} X X p{2.5cm} @{}}
\caption{RGAIG+ Consumer vs Clinical EEG System Comparison}
\label{tab:rgaig_consumer_vs_clinical}\\
\toprule
\thead{Metric} & \thead{Consumer Wearable EEG} & \thead{Clinical EEG} & \thead{Trade-off} \\
\midrule
\endfirsthead
\endhead
\bottomrule
\endlastfoot
Channels        & 4--16 (10--20 subset)                        & 32--128 (full 10--20 or 10--10)               & Lower spatial resolution; sufficient for screening \\
\addlinespace
Signal quality  & Moderate (dry/semi-dry electrodes)            & High (wet gel electrodes)                     & Higher noise floor; AI compensates via filtering \\
\addlinespace
Usability       & High (self-applied, $<$5 min setup)          & Low (technician-applied, 30--60 min)          & Consumer access vs signal fidelity \\
\addlinespace
Mobility        & Fully portable; home/ambulatory monitoring    & Stationary; hospital/clinic setting            & Continuous monitoring vs controlled environment \\
\addlinespace
Cost            & \$300--\$1{,}000 (Emotiv, Muse, OpenBCI)     & \$15{,}000--\$100{,}000 (clinical systems)    & Democratised access vs gold-standard precision \\
\addlinespace
Recording       & Hours--days (ambulatory)                      & 30 min--72 hours (inpatient)                  & Longitudinal data vs snapshot \\
\addlinespace
Clinical use    & Screening, monitoring, trend detection        & Diagnosis, presurgical mapping, research      & Complementary roles; not replacement \\
\addlinespace
AI suitability  & Optimised for DL pattern recognition       & Full-spectrum analysis; visual interpretation  & AI compensates reduced channels with learned features \\
\addlinespace
\midrule
\multicolumn{4}{@{}p{\dimexpr\textwidth-2\tabcolsep}@{}}{\textit{RGAIG+ positions consumer EEG as complementary screening tool, not clinical replacement. AI models trained on consumer-grade signals; accuracy validated: ASD 97.67\%, Epilepsy 99.02\%. Clinician escalation mandatory for positive screening results (governance Layer~7).}} \\
\end{xltabular}
\endgroup

\vspace{1.5em}

%% ── Table 5: Wearable Device Validation Framework ──
Table~\ref{tab:rgaig_device_validation} documents the six-level validation framework for wearable EEG devices, specifying acceptance criteria and applicable standards.

\needspace{4\baselineskip}
\begingroup\singlespacing\tablefontsize\tablestripes
\begin{xltabular}{\textwidth}{@{} p{2.2cm} X X p{1.8cm} @{}}
\caption{RGAIG+ Wearable Device Validation Framework}
\label{tab:rgaig_device_validation}\\
\toprule
\thead{Validation Type} & \thead{Description} & \thead{Acceptance Criteria} & \thead{Standard} \\
\midrule
\endfirsthead
\endhead
\bottomrule
\endlastfoot
Signal validation       & Compare wearable EEG output with clinical-grade EEG on same subjects; assess frequency band correlation & Pearson $r \geq 0.80$ for alpha, beta, theta bands; SNR $\geq$ 10~dB across all channels & IEC 60601 \\
\addlinespace
Algorithm validation    & Verify AI diagnostic predictions against ground-truth clinical labels; assess generalisation across devices & Accuracy $\geq$ 85\%; AUC $\geq$ 0.90; F1 $\geq$ 0.85 per domain; LOSO validation & FDA SaMD \\
\addlinespace
Clinical validation     & Neurologist evaluation of AI-flagged cases; sensitivity/specificity against clinical diagnosis & Sensitivity $\geq$ 90\%; specificity $\geq$ 85\%; expert agreement $\kappa \geq$ 0.70 & IMDRF \\
\addlinespace
Usability validation    & End-user testing with target population; task completion, error rates, and satisfaction measurement & SUS $\geq$ 70; task completion $\geq$ 90\%; critical use errors = 0 & IEC 62366-1 \\
\addlinespace
Safety validation       & Electrical safety testing; skin biocompatibility; thermal assessment; electromagnetic compatibility & IEC 60601 compliance; no adverse skin reactions in 48-hour patch test; EMC within limits & ISO 14971 \\
\addlinespace
Reliability validation  & Test-retest reliability across sessions; inter-device reproducibility; long-term stability assessment & ICC $\geq$ 0.90 test-retest; inter-device CV $<$ 10\%; 30-day stability within $\pm$5\% & ISO 13485 \\
\addlinespace
\midrule
\multicolumn{4}{@{}p{\dimexpr\textwidth-2\tabcolsep}@{}}{\textit{Validation protocol follows V-model approach: unit $\rightarrow$ integration $\rightarrow$ system $\rightarrow$ clinical validation. Each validation stage must be completed before proceeding to next RGAIG+ maturity level.}} \\
\end{xltabular}
\endgroup

\vspace{1.5em}

%% ── Table 6: Device Reliability Metrics ──
Table~\ref{tab:rgaig_device_reliability} reports the eight device reliability and monitoring metrics tracked by the RGAIG+ monitoring layer.

\needspace{4\baselineskip}
\begingroup\singlespacing\tablefontsize\tablestripes
\begin{xltabular}{\textwidth}{@{} p{2.8cm} X p{1.8cm} p{2.2cm} @{}}
\caption{RGAIG+ Device Reliability and Monitoring Metrics}
\label{tab:rgaig_device_reliability}\\
\toprule
\thead{Metric} & \thead{Description} & \thead{Target} & \thead{Monitoring} \\
\midrule
\endfirsthead
\endhead
\bottomrule
\endlastfoot
Signal-to-noise ratio   & EEG signal quality relative to background noise; measured per channel per session & $\geq$ 10~dB & Real-time \\
\addlinespace
Electrode impedance     & Contact quality between electrode and scalp; indicator of signal reliability & $<$ 20~k$\Omega$ & Per-session \\
\addlinespace
Data loss rate          & Percentage of EEG samples lost during Bluetooth/Wi-Fi transmission & $<$ 1\% & Real-time \\
\addlinespace
Battery endurance       & Continuous recording time per full charge under diagnostic workload & $\geq$ 6 hours & Per-charge \\
\addlinespace
Calibration drift       & Signal baseline shift between recording sessions; affects longitudinal comparison & $<$ 5\% & Weekly \\
\addlinespace
Device uptime           & Percentage of scheduled recording time with valid signal acquisition & $\geq$ 95\% & Daily \\
\addlinespace
Reconnection time       & Time to re-establish connection after dropout event & $<$ 5 seconds & Per-event \\
\addlinespace
Thermal safety          & Device surface temperature during extended recording (skin contact area) & $<$ 41$^{\circ}$C & Continuous \\
\addlinespace
\midrule
\multicolumn{4}{@{}p{\dimexpr\textwidth-2\tabcolsep}@{}}{\textit{Metrics tracked by RGAIG+ Monitoring Layer (Layer~8). Threshold violations trigger automated alerts and governance escalation per KPI protocol (Table~\ref{tab:rgaig_device_risk}).}} \\
\end{xltabular}
\endgroup

\vspace{1.5em}

%% ── Table 7: Device Safety Evaluation ──
Table~\ref{tab:rgaig_device_safety} presents the seven safety aspects evaluated for consumer wearable EEG devices, with associated risks and mitigations.

\needspace{4\baselineskip}
\begingroup\singlespacing\tablefontsize\tablestripes
\begin{xltabular}{\textwidth}{@{} p{2.8cm} X X @{}}
\caption{RGAIG+ Wearable Device Safety Evaluation}
\label{tab:rgaig_device_safety}\\
\toprule
\thead{Safety Aspect} & \thead{Risk Description} & \thead{Mitigation and Standard} \\
\midrule
\endfirsthead
\endhead
\bottomrule
\endlastfoot
Electrical safety       & Risk of excessive current through scalp electrodes; potential for burns or tissue damage & Current leakage $<$ 10~$\mu$A per IEC 60601-1; isolated power supply; automatic shutdown on fault detection \\
\addlinespace
Skin biocompatibility   & Extended electrode contact may cause irritation, allergic reaction, or pressure sores & ISO 10993 biocompatibility testing; hypoallergenic electrode materials; recommended wear duration limits \\
\addlinespace
Thermal safety          & Device heating during prolonged recording; battery thermal runaway risk & Surface temperature monitoring ($<$ 41$^{\circ}$C); thermal shutdown circuit; battery protection IC \\
\addlinespace
Electromagnetic safety  & EMI from device affecting other medical equipment; susceptibility to external EMI & IEC 60601-1-2 EMC compliance; shielded electronics; frequency band isolation \\
\addlinespace
Medical guidance safety & AI-generated diagnostic suggestions may be incorrect or misleading for consumers & Mandatory confidence thresholds ($\geq$ 0.85); disclaimer on all outputs; clinician escalation for positive screens \\
\addlinespace
Anomaly detection       & Device malfunction producing artefactual signals misinterpreted as pathological & Real-time signal quality monitoring; impedance checks; automatic flagging of artefact-contaminated epochs \\
\addlinespace
Data security           & Neurological data interception during wireless transmission; unauthorised access & AES-256 encryption; secure BLE pairing; HIPAA/GDPR-compliant data handling per ISO/IEC 27001 \\
\addlinespace
\midrule
\multicolumn{3}{@{}p{\dimexpr\textwidth-2\tabcolsep}@{}}{\textit{Safety evaluation follows ISO 14971 risk management process: hazard identification $\rightarrow$ risk estimation $\rightarrow$ risk evaluation $\rightarrow$ risk control $\rightarrow$ residual risk assessment. All mitigations traceable to standards.}} \\
\end{xltabular}
\endgroup

\vspace{1.5em}

%% ── Table 8: Wearable Device Testing Framework ──
Table~\ref{tab:rgaig_device_testing} summarises the six-level testing framework for wearable EEG devices, from hardware through security testing.

\needspace{4\baselineskip}
\begingroup\singlespacing\tablefontsize\tablestripes
\begin{xltabular}{\textwidth}{@{} p{2.8cm} X X @{}}
\caption{RGAIG+ Wearable Device Testing Framework}
\label{tab:rgaig_device_testing}\\
\toprule
\thead{Testing Level} & \thead{Scope and Description} & \thead{Example Tests} \\
\midrule
\endfirsthead
\endhead
\bottomrule
\endlastfoot
Hardware testing        & Physical device components: electrodes, circuit board, battery, enclosure & Electrode impedance stability; drop test (1.5~m); IP rating verification; battery cycle testing (500 cycles) \\
\addlinespace
Signal testing          & EEG signal acquisition and quality under various conditions & Noise floor measurement; frequency response (0.5--100~Hz); common mode rejection ratio; motion artefact assessment \\
\addlinespace
Software testing        & Device firmware and companion app functionality & Firmware update reliability; data synchronisation; error handling; memory leak testing; edge case inputs \\
\addlinespace
Integration testing     & End-to-end device-to-AI pipeline connectivity & Bluetooth pairing; data streaming continuity; AI model inference trigger; result display latency; failover handling \\
\addlinespace
Stress testing          & Performance under sustained or adverse conditions & 24-hour continuous recording; multi-user concurrent streaming; low-battery operation; weak connectivity environments \\
\addlinespace
Security testing        & Data protection and access control verification & Penetration testing on BLE connection; encryption validation; authentication bypass attempts; data-at-rest protection \\
\addlinespace
\midrule
\multicolumn{3}{@{}p{\dimexpr\textwidth-2\tabcolsep}@{}}{\textit{Testing follows V-model: hardware $\rightarrow$ signal $\rightarrow$ software $\rightarrow$ integration $\rightarrow$ stress $\rightarrow$ security. All test results documented in device validation report; traceability matrix links tests to requirements.}} \\
\end{xltabular}
\endgroup

\vspace{1.5em}

%% ── Table 9: Consumer Trust Metrics ──
Table~\ref{tab:rgaig_consumer_trust_metrics} reports the consumer trust and adoption metrics used to evaluate RGAIG+ deployment outcomes.

\needspace{4\baselineskip}
\begingroup\singlespacing\tablefontsize\tablestripes
\begin{xltabular}{\textwidth}{@{} p{2.8cm} X p{2.8cm} p{1.8cm} @{}}
\caption{RGAIG+ Consumer Trust and Adoption Metrics}
\label{tab:rgaig_consumer_trust_metrics}\\
\toprule
\thead{Metric Category} & \thead{Description} & \thead{Measurement} & \thead{Target} \\
\midrule
\endfirsthead
\endhead
\bottomrule
\endlastfoot
\multicolumn{4}{@{}p{\dimexpr\textwidth-2\tabcolsep}@{}}{\textbf{\color{ggunavy}Consumer Trust Metrics}} \\
\midrule
Satisfaction score      & Overall user satisfaction with diagnostic service & Likert 1--5 survey & $\geq$ 4.0/5 \\
\addlinespace
Recommendation rate     & Willingness to recommend service to others (NPS proxy) & Binary yes/no + NPS scale & NPS $\geq$ 50 \\
\addlinespace
Retention rate          & Percentage of users continuing beyond 90 days & Usage tracking & $\geq$ 70\% \\
\addlinespace
Clinical follow-up rate & Percentage of flagged users consulting clinicians & Follow-up survey & $\geq$ 80\% \\
\addlinespace
Explanation clarity     & Perceived clarity of RAG-generated diagnostic explanations & Likert 1--5 & $\geq$ 3.5/5 \\
\addlinespace
\midrule
\multicolumn{4}{@{}p{\dimexpr\textwidth-2\tabcolsep}@{}}{\textbf{\color{ggunavy}Consumer Adoption Metrics (TAM-TOE-RBV)}} \\
\midrule
Perceived usefulness    & Belief that the system improves health monitoring outcomes & TAM PU scale (4 items) & $\geq$ 4.0/5 \\
\addlinespace
Perceived ease of use   & Belief that the system is easy to learn and operate & TAM PEOU scale (3 items) & $\geq$ 4.0/5 \\
\addlinespace
Behavioural intention   & Intention to continue using the system & TAM BI scale (2 items) & $\geq$ 4.0/5 \\
\addlinespace
System usability        & Objective usability assessment of consumer interface & SUS questionnaire & $\geq$ 70/100 \\
\addlinespace
Comfort rating          & Physical comfort during extended wearable use & Likert 1--5 & $\geq$ 4.0/5 \\
\addlinespace
\midrule
\multicolumn{4}{@{}p{\dimexpr\textwidth-2\tabcolsep}@{}}{\textit{Trust metrics adapted from McKnight et al.\ (2002) and Siau \& Wang (2018). Adoption metrics from \citet{davis1989tam} TAM. Metrics collected at 30, 60, and 90-day intervals post-deployment.}} \\
\end{xltabular}
\endgroup

\vspace{1.5em}

%% ── Table 10: Responsible Device Governance (RACI Matrix) ──
Table~\ref{tab:rgaig_device_governance} maps the seven stakeholder roles to their responsibilities and accountability domains within the RGAIG+ governance model.

\needspace{4\baselineskip}
\begingroup\singlespacing\tablefontsize\tablestripes
\begin{xltabular}{\textwidth}{@{} p{2.8cm} X p{2.8cm} @{}}
\caption{RGAIG+ Responsible Device Governance --- Stakeholder Accountability}
\label{tab:rgaig_device_governance}\\
\toprule
\thead{Stakeholder} & \thead{Responsibility} & \thead{Account\-ability} \\
\midrule
\endfirsthead
\endhead
\bottomrule
\endlastfoot
Device manufacturer     & Hardware quality, electrode safety, firmware reliability, IEC 60601 compliance, and post-market surveillance & Device safety \\
\addlinespace
AI developer            & Model accuracy, fairness testing, hallucination control, bias monitoring, and model versioning & Model reliability \\
\addlinespace
Healthcare provider     & Clinical validation, output interpretation, patient communication, and escalation response & Clinical oversight \\
\addlinespace
Governance board        & Policy enforcement, risk threshold approval, audit oversight, and regulatory liaison & Policy compliance \\
\addlinespace
Data controller         & Data privacy, HIPAA/GDPR compliance, consent management, de-identification, and breach notification & Data protection \\
\addlinespace
Consumer                & Correct device usage, self-reporting of issues, informed consent, and clinical follow-up on positive screens & Self-care adherence \\
\addlinespace
Regulatory authority    & Standards enforcement, market authorisation, adverse event reporting, and periodic audit & Legal compliance \\
\addlinespace
\midrule
\multicolumn{3}{@{}p{\dimexpr\textwidth-2\tabcolsep}@{}}{\textit{Governance model follows RACI (Responsible, Accountable, Consulted, Informed) principles. All stakeholders documented in RGAIG+ governance charter with escalation pathways.}} \\
\end{xltabular}
\endgroup

\vspace{1.5em}

%% ── Table 11: Ethical Considerations for Consumer EEG Devices ──
Table~\ref{tab:rgaig_device_ethics} documents the six ethical considerations for consumer EEG diagnostics, with corresponding RGAIG+ mitigations.

\needspace{4\baselineskip}
\begingroup\singlespacing\tablefontsize\tablestripes
\begin{xltabular}{\textwidth}{@{} p{2.8cm} X X @{}}
\caption{RGAIG+ Ethical Considerations for Consumer EEG Diagnostics}
\label{tab:rgaig_device_ethics}\\
\toprule
\thead{Ethical Issue} & \thead{Risk Description} & \thead{RGAIG+ Mitigation} \\
\midrule
\endfirsthead
\endhead
\bottomrule
\endlastfoot
Neurological privacy    & EEG data reveals cognitive states, emotional responses, and neurological conditions --- highly sensitive personal data & End-to-end encryption (AES-256); de-identification at source; GDPR Article 9 special category processing; data minimisation \\
\addlinespace
Misinterpretation       & Consumer may misunderstand AI-generated diagnostic screening results as definitive clinical diagnosis & Mandatory disclaimer on all outputs; confidence score displayed; language: ``screening suggestion'' not ``diagnosis''; clinician escalation \\
\addlinespace
Over-reliance on AI     & Consumer may substitute AI screening for professional medical consultation & Clear disclaimers; ``This is not medical advice'' on all screens; automated prompts to consult healthcare provider \\
\addlinespace
Algorithmic bias        & AI models may perform differently across demographic groups (age, sex, ethnicity) & Fairness metrics per subgroup; bias detection in monitoring layer; demographic parity and equalised odds thresholds \\
\addlinespace
Informed consent        & Consumer must understand what data is collected, how it is used, and the limitations of AI screening & Plain-language consent form; opt-in data sharing; granular consent per use case; right to withdraw and data deletion \\
\addlinespace
Digital divide          & Consumer EEG systems may be inaccessible to underserved populations, widening health disparities & Cost-accessible device options (\$300--\$500); multi-language interface; low-bandwidth compatibility; community health worker support \\
\addlinespace
\midrule
\multicolumn{3}{@{}p{\dimexpr\textwidth-2\tabcolsep}@{}}{\textit{Ethical framework aligned with WHO AI for Health (2021) six guiding principles and OECD AI Principles. Ethics review conducted as part of RGAIG+ Governance Layer (Layer~9) at each maturity level transition.}} \\
\end{xltabular}
\endgroup

\clearpage
